% Píše: aplikace. Závěr shrnuje obě části projektu.
\chapter{Závěr}
\label{chap:zaver}

Projekt splnil variantu~2 zadání, komunikaci s~okolím pomocí bezdrátové komunikace, a~to
v~komplexní podobě přes Bluetooth Low Energy. Vznikl chytrý dům řízený Raspberry Pi 3B+, který
měří skutečná čidla, simuluje vnitřní teplotu s~termostatem, automaticky řídí barevné světlo
a~vystupuje jako periferie BLE, a~mobilní aplikace pro Android, která se k~němu připojí jako
klient. Oba příklady, které zadání u~komplexní komunikace uvádí, projekt pokrývá: aplikace sbírá
data ze senzorů přes Bluetooth LE a~vizualizuje řídicí technologii, tedy termostat a~světlo,
v~přehledu, grafu teplot a~deníku událostí. Navíc dům i~ovládá: mění nastavení a~posílá příkazy.
Stav přichází každou sekundu a~změnu nastavení dům potvrdí notifikací, takže aplikace ukazuje
skutečný stav bez prodlevy.

Obě části spojuje přesně popsané komunikační rozhraní (kapitola~\ref{chap:rozhrani}). Díky němu
se mohly vyvíjet a~testovat souběžně a~samostatně: backend proti falešnému hardwaru
a~v~režimu simulace na PC (oddíl~\ref{sec:simulace}), aplikace proti simulovanému domu, který
dodržuje stejná pravidla (oddíl~\ref{sec:app-simulace}). Backend má 746 a~aplikace 1\,173
automatických testů a~všechny procházejí. Součástky se vyzkoušely na Raspberry Pi jednotlivě
i~v~běžící službě a~společný test ověřil připojení s~MTU 517, notifikace a~zápisy z~tabletu, které
dům uložil. Problémy, které se objevily až na hardwaru, se opravily (oddíl~\ref{sec:problemy-hw}).

Řešení má tato omezení:
\begin{itemize}
  \item Vnitřní teplota je simulovaná a~topení, chlazení i~světlo zastupují LED. Dům nemá
        fyzický model ani výkonové výstupy.
  \item Čidlo BMP280 neměří vlhkost, a~vlhkost proto v~sestavě chybí.
  \item Rozhraní nemá zabezpečení: připojit se a~měnit nastavení může kdokoli v~dosahu
        (oddíl~\ref{sec:zabezpeceni}).
  \item Raspberry Pi nemá hodiny reálného času. Bez připojení k~internetu musí čas nastavit tablet.
  \item Aplikace drží historii grafu a~deník jen v~paměti, nejvýš 6~hodin a~200~událostí, a~po
        zavření o~ně přijde.
  \item Výsledky některých ručních zkoušek nejsou zaznamenané (tabulka~\ref{tab:hw-vysledky}).
\end{itemize}

Projekt jde rozšířit o~šifrované a~ověřené zápisy se spárováním tabletu, o~ukládání historie
měření na Raspberry Pi a~její přenos do aplikace, o~skutečné akční členy (relé pro topení
a~ventilátor, světelný pásek s~vlastním napájením), o~čidlo BME280 s~vlhkostí, kterou už obě
části umějí zobrazit, a~o~hodiny reálného času. V~aplikaci by se hodilo upozornění na poruchu
čidla nebo ztrátu spojení i~ve chvíli, kdy aplikace neběží na popředí.
